\section{Описание}
Требуется найти лексикографически минимальный циклический сдвиг строки $S$ длины $L$.

Как описано в \cite{Gasfield}, суффиксное дерево является мощнейшим индексом для решения широкого класса строковых задач. Для того чтобы получить доступ ко всем циклическим сдвигам строки, достаточно удвоить исходную строку: $S' = S + S$. Теперь каждый циклический сдвиг $S$ является подстрокой длины $L$ в строке $S'$.

Для решения задачи к строке $S'$ добавляется уникальный терминальный символ (например, \texttt{\#}), после чего за время $O(N)$ строится суффиксное дерево с помощью алгоритма Укконена. 
Поиск минимального сдвига сводится к жадному спуску от корня дерева: на каждом шаге выбирается ребро, начинающееся с лексикографически минимального символа (исключая \texttt{\#}). Спуск продолжается, пока не будет накоплено $L$ символов.

\textbf{Особенности реализации:}
Алгоритм Укконена требует сложного управления памятью. В данной реализации применены следующие оптимизации:
\begin{enumerate}
\item Вместо динамического выделения памяти (\texttt{new Node}) узлы хранятся в непрерывном массиве \texttt{std::vector<Node>}. Это исключает утечки памяти и ускоряет аллокацию.
\item Переходы по ребрам реализованы через \texttt{std::map<char, int>}. Поскольку красно-черное дерево внутри \texttt{map} автоматически сортирует ключи по возрастанию, при поиске минимального сдвига достаточно всегда брать первого потомка (игнорируя \texttt{\#}).
\end{enumerate}

\section{Исходный код}
Основная структура узла дерева.

\begin{lstlisting}[language=C++]
struct Node {
    int start;
    int end;
    int link;
    map<char, int> next;

    Node(int s, int e, int l) : start(s), end(e), link(l) {}

    int length(int leafEnd) {
        if (start == -1) return 0;
        return (end == INF ? leafEnd : end) - start + 1;
    }
};
\end{lstlisting}

Метод спуска по суффиксному дереву для поиска ответа. Использование \texttt{cout.write} позволяет избежать медленного копирования строк.

\begin{lstlisting}[language=C++]
void printSmallest(int L) {
    int curr = root;
    while (L > 0) {
        int best_node = -1;
        for (auto const& [ch, node_idx] : tree[curr].next) {
            if (ch == '#') continue;
            best_node = node_idx;
            break; 
        }
        if (best_node == -1) break;

        int edge_len = tree[best_node].length(leafEnd);
        int chars_to_take = min(edge_len, L);
        
        cout.write(&text[tree[best_node].start], chars_to_take);
        
        L -= chars_to_take;
        curr = best_node;
    }
    cout << "\n";
}
\end{lstlisting}

Полный листинг программы (включая логику алгоритма Укконена):

\lstinputlisting[language=C++, caption={Линеаризация строки через суффиксное дерево}, label={lst:maincpp}]{../main.cpp}
\pagebreak

\section{Консоль}
Пример сборки программы и тестирования.

\begin{alltt}
$ g++ -std=c++20 -Wall -Wextra -O3 main.cpp -o tree_build.out
$ cat test.in
xabcd
$ ./tree_build.out < test.in > result.out
$ cat result.out
abcdx
\end{alltt}
\pagebreak